\SGAchapter{II}{基本群的有限性性质}{Propriétés de finitude du groupe fondamental}
\SourcePage{34}
\begin{center}Michèle Raynaud\end{center}
设 $k$ 是\Term{characteristic}{特征} $p\geq0$ 的\Term{separably-closed-field}{可分闭域}，$X$ 是 $k$ 上\Term{connected}{连通}的\Term{scheme}{概形}，且为\Term{finite-type}{有限型}；以 $\Pi_1^{(p)}(X)$ 表示 $\Pi_1(X)$ 的最大“与 $p$ 互素”商。本讲证明：若 $X$ 是\Term{surface}{曲面}，则 $\Pi_1^{(p)}(X)$ \Term{topologically-finitely-presented}{在拓扑意义下具有有限表示}；若承认\Term{resolution-singularities}{奇点消解}，则任意 $X$ 也有此性质。

在\Term{surface}{曲面}的情形，由 J.-P. Murre 提出的方法是归约到 $\mathbf P^1$ 上的一个\Term{fibration}{纤维化}，它具有下述命题的性质 (i)、(ii)、(iii)。
\begin{SGAresult}{命题}{2.1}
设 $k$ 是\Term{algebraically-closed-field}{代数闭域}，$S$ 是 $k$ 上\Term{projective}{射影}、\Term{irreducible}{不可约}、\Term{smooth}{光滑}的\Term{surface}{曲面}。则存在一个\Term{blow-up}{爆破}
\[g:S'\longrightarrow S,\]
其中 $S'$ 是\Term{regular}{正则}\Term{surface}{曲面}，且带有\Term{fibration}{纤维化}
\[f:S'\longrightarrow\mathbf P^1,\]
使 $f$ 有一个\Term{section}{截面} $\sigma$，并满足：
\begin{enumerate}
\item[(i)] $f$ 在 $\sigma(\mathbf P^1)$ 的各点\Term{smooth}{光滑}；
\item[(ii)] $f$ 的\Term{fiber}{纤维}是\Term{geometrically-integral}{几何整}\Term{curve}{曲线}，且只有普通\Term{singular-point}{奇点}；
\item[(iii)] 存在 $\mathbf P^1$ 的\Term{nonempty}{非空}\Term{open-subset}{开集} $U$，使 $f$ 在 $U$ 中各点处的\Term{fiber}{纤维}\Term{smooth}{光滑}。
\end{enumerate}
\end{SGAresult}
\begin{proof}
存在\Term{projective-embedding}{射影嵌入} $S\hookrightarrow\mathbf P^r$，使相对于这个\Term{embedding}{嵌入}可以找到一个 \Term{lefschetz-pencil}{Lefschetz 铅笔}（第 XVII 讲 2.5）；命题由此\Term{pushout}{推出}。事实上，设 $D=\{H_t\}_{t\in\mathbf P^1}$ 是这样一个\Term{hyperplane-pencil}{超平面铅笔}。回顾其下述性质（同上，2.2）。
\SourcePage{35}
\begin{enumerate}
\item[a)] $D$ 的轴 $\Delta$ 与 $S$ \Term{transverse-intersection}{横截相交}；因此\Term{closed-subscheme}{闭子概形} $\Delta\cap S$ 是\Term{reduced}{约化}的，由\Term{finite}{有限}个\Term{closed-point}{闭点} $x_1,\ldots,x_n$ 构成。
\item[b)] 存在 $\mathbf P^1$ 的\Term{nonempty}{非空}\Term{open-subset}{开集} $U$，使任意 $t\in U$ 对应的 $H_t$ 与 $S$ \Term{transverse-intersection}{横截相交}。
\item[c)] 对每一点 $t\in\mathbf P^1-U$，$H_t$ 除一点之外都与 $S$ \Term{transverse-intersection}{横截相交}，而这一点是\Term{ordinary-quadratic-singularity}{普通二次奇点}。
\end{enumerate}
对于 $S$ 中不属于 $\Delta$ 的每一点 $x$，恰有一个\Term{hyperplane}{超平面} $H_t$ 经过它。把 $x$ 映到 $t\in\mathbf P^1$ 所得的映射，是\Term{domain-of-definition}{定义域}为 $S-(\Delta\cap S)$ 的\Term{rational-map}{有理映射} $a:S\dashrightarrow\mathbf P^1$。

考虑 $S$ 沿\Term{closed-subscheme}{闭子概形} $S\cap\Delta$ 的\Term{blow-up}{爆破} $g:S'\to S$。\Term{scheme}{概形} $S'$ \Term{regular}{正则}，点 $x_i$ 上方的\Term{fiber}{纤维} $S'_{x_i}$ 等同于由 $S$ 在 $x_i$ 处的\Term{tangent-space}{切空间}定义的\Term{projective-bundle}{射影丛}，所以\Term{isomorphism}{同构}于 $\mathbf P^1$（EGA IV 19.4.2、19.4.4）。此外，\Term{rational-map}{有理映射} $f=ag$ 处处有定义（第 XVIII 讲 3.1）。最后，每个 $x_i$ 给出 $f$ 的一个\Term{section}{截面}：它把 $t\in\mathbf P^1$ 映到\Term{curve}{曲线} $H_t\cap S$ 在 $x_i$ 处的切向量。

验证 (i)、(ii)、(iii)。注意，对每个 $t\in\mathbf P^1$，\Term{curve}{曲线} $S'_t$ 与 $S_t=H_t\cap S$ \Term{isomorphism}{同构}：\Term{morphism}{态射} $S'_t\to S_t$ 除可能在 $x_i$ 处之外是\Term{isomorphism}{同构}，而 $S_t$ 在这些点\Term{regular}{正则}；又因 $S'_t$ 是 $S'$ 中的\Term{complete-intersection}{完全交}，故 $S'_t$ 是整的，于是\Term{isomorphism}{同构}于 $S_t$。

条件 (i)、(ii) 正是条件 a)、b)。最后，$f$ 在 $\sigma$ 的各点\Term{smooth}{光滑}，因为 $f$ \Term{flat}{平坦}，且\Term{fiber}{纤维} $S'_t$ 在 $\sigma(t)$ 处\Term{geometrically-regular}{几何正则}。
\end{proof}
\SourcePage{36}
\begin{SGAresult}{推论}{2.2}
沿用 2.1 的记号，设 $Y$ 是 $S$ 中的\Term{normal-crossings-divisor}{正规交叉除子}，$Y'=Y\times_S S'$。则可选择 $f,g,U$，使 $Y'\times_{\mathbf P^1}U$ 是相对于 $U$ 的\Term{normal-crossings-divisor}{正规交叉除子}。
\end{SGAresult}
\begin{proof}
已知，若某个\Term{embedding}{嵌入} $S\hookrightarrow\mathbf P^r$ 存在 \Term{lefschetz-pencil}{Lefschetz 铅笔}，则这些\Term{pencil}{铅笔}构成 $\mathbf P^r$ 中直线簇的一个\Term{open-subset}{开集}（第 XVII 讲 3.2.1）。下述 \Term{lefschetz-pencil}{Lefschetz 铅笔}的\Term{set}{集合}也如此：其轴不与 $Y$ 相交，且存在 $\mathbf P^1$ 的\Term{nonempty}{非空}\Term{open-subset}{开集} $U_1$，使 $t\in U_1$ 时 $Y\cap H_t$ \Term{smooth}{光滑}。因此只须从满足这两个条件的\Term{pencil}{铅笔}构造 $f,g$，再把 $U$ 换成 $U\cap U_1$。
\end{proof}
\textbf{2.3.0.} 设 $G$ 是\Term{profinite-group}{投有限群}。回顾：若存在整数 $n$，使 $G$ \Term{isomorphism}{同构}于 $n$ 个\Term{generator}{生成元}的自由投群 $F(n)$ 的一个商，则称 $G$ \Term{topologically-finitely-generated}{在拓扑意义下有限生成}。令 $K=\operatorname{Ker}(F(n)\to G)$；若此外能取\Term{finite}{有限}个元素 $k_1,\ldots,k_r\in K$，使 $K$ 中包含这些元素的最小闭\Term{normal-subgroup}{正规子群}等于 $K$，则称 $G$ \Term{topologically-finitely-presented}{在拓扑意义下具有有限表示}。
\begin{SGAresult}{\Term{theorem}{定理}}{2.3.1}
设 $k$ 是\Term{characteristic}{特征} $p\geq0$ 的\Term{separably-closed-field}{可分闭域}，$X$ 是 $k$ 上\Term{connected}{连通}的\Term{finite-type}{有限型}\Term{scheme}{概形}。假设 $k$ 的一个\Term{algebraic-closure}{代数闭包}上，\Term{dimension}{维数} $\leq\dim(X)$ 的\Term{finite-type}{有限型}\Term{scheme}{概形}均可强\Term{resolution}{消解}\Term{singular-point}{奇点}（SGA 5 I 3.1.5；当 $\dim X=2$ 时由 [1] 知此条件成立；当 $k=0$ 时由 [2] 知其成立）。则\Term{fundamental-group}{基本群} $\Pi_1^{(p)}(X)$ \Term{topologically-finitely-presented}{在拓扑意义下具有有限表示}。
\end{SGAresult}
\begin{proof}
由 SGA 1 IX 4.10，可假设 $k$ \Term{algebraically-closed}{代数闭}。

\textbf{1) 归约到 $X$ 为\Term{regular}{正则}\Term{surface}{曲面}的情形。}
若 $f:X'\to X$ 是\Term{finite-type}{有限型}\Term{morphism}{态射}，令 $X''=X'\times_X X'$，并分别以 $(X'_a)_{a\in A}$、$(X''_b)_{b\in B}$ 表示 $X'$、$X''$ 的\Term{connected-component}{连通分支}。当 $f$ 是\Term{etale-covering}{平展覆盖}\Term{category}{范畴}的\Term{effective-descent-morphism}{有效下降态射}时，\Term{fundamental-group}{基本群} $\Pi_1(X)$ 可用 $\Pi_1(X'_a)$、$\Pi_1(X''_b)$ 显式计算。
\SourcePage{37}
由此（SGA 1 IX 5.2、5.3），若各 $\Pi_1(X'_a)$ \Term{topologically-finitely-generated}{在拓扑意义下有限生成}，则 $\Pi_1(X)$ 也如此；若各 $\Pi_1(X'_a)$ \Term{topologically-finitely-presented}{在拓扑意义下具有有限表示}，且各 $\Pi_1(X''_b)$ \Term{topologically-finitely-generated}{在拓扑意义下有限生成}，则 $\Pi_1(X)$ \Term{topologically-finitely-presented}{在拓扑意义下具有有限表示}。

因 $X$ 可\Term{resolution}{消解}\Term{singular-point}{奇点}，可取\Term{regular-scheme}{正则概形} $X'$ 和\Term{proper}{固有}\Term{birational-morphism}{双有理态射} $f:X'\to X$。由于 $f$ 是\Term{etale-covering}{平展覆盖}\Term{category}{范畴}的\Term{effective-descent-morphism}{有效下降态射}，只须在附加假设 $X$ \Term{regular}{正则}之下证明\Term{theorem}{定理}。这样首先\Term{pushout}{推出} $\Pi_1(X)$ \Term{topologically-finitely-generated}{在拓扑意义下有限生成}；知道 $\Pi_1(X')$ \Term{topologically-finitely-presented}{在拓扑意义下具有有限表示}，且各 $\Pi_1(X''_b)$ \Term{topologically-finitely-generated}{在拓扑意义下有限生成}后，就\Term{pushout}{推出} $\Pi_1(X)$ \Term{topologically-finitely-presented}{在拓扑意义下具有有限表示}。

现在，设 $(U_i)_{i\in I}$ 是 $X$ 的\Term{finite}{有限}\Term{affine-open-cover}{仿射开覆盖}，$f:\coprod U_i\to X$ 是\Term{canonical-morphism}{典范态射}。因 $f$ 是\Term{etale-covering}{平展覆盖}\Term{category}{范畴}的\Term{effective-descent-morphism}{有效下降态射}，只须对各 $U_i$ 证明\Term{theorem}{定理}；因此归约到 $X$ \Term{affine}{仿射}且\Term{smooth}{光滑}的情形。

以下设 $X$ 是 $k$ 上的\Term{smooth}{光滑}\Term{affine-scheme}{仿射概形}，$\dim X>2$。此时可应用\Term{quasi-projective-scheme}{拟射影概形}的 Lefschetz \Term{theorem}{定理}（[3] V 7.4）：若 $Y$ 是 $X$ 的一个“充分\Term{generic}{一般}”的\Term{hyperplane-section}{超平面截面}，则有\Term{isomorphism}{同构}
\[\Pi_1(Y)\simeq\Pi_1(X).\]
故只须对 $Y$ 证明\Term{theorem}{定理}，这便归约到 $\dim X=2$ 的情形。

\textbf{2) $X$ 为 $k$ 上\Term{smooth}{光滑}\Term{surface}{曲面}的情形。}
由\Term{surface}{曲面}的强\Term{resolution-singularities}{奇点消解}\Term{theorem}{定理}，可取\Term{smooth}{光滑}、整、\Term{projective-surface}{射影曲面} $S$ 及\Term{open-immersion}{开浸入} $i:X\hookrightarrow S$，使\Term{divisor}{除子} $Y=S-X$ \Term{normal-crossings}{正规交叉}。应用推论 2.2，并沿用其中记号。
\SourcePage{38}
说明只须证明 $X'=X\times_S S'$ 的\Term{fundamental-group}{基本群}\Term{topologically-finitely-presented}{在拓扑意义下具有有限表示}。因为 $X,X'$ \Term{regular}{正则}，且
\[\operatorname{codim}\bigl(X-\!\bigcup_{1\leq i\leq n}\{x_i\},X\bigr)\geq2,\]
$X'$ 的任意\Term{etale-covering}{平展覆盖}在 $X-\bigcup_{1\leq i\leq n}\{x_i\}$ 上诱导一个\Term{etale-covering}{平展覆盖}，而它唯一延拓为 $X$ 的\Term{etale-covering}{平展覆盖}（SGA 2 1.11）；于是有\Term{isomorphism}{同构}
\[\Pi_1(X')\simeq\Pi_1(X).\]
考虑\Term{cartesian-diagram}{笛卡儿图}
\[\begin{tikzcd}
X'\arrow[d,"h"]&X'_U\arrow[l]\arrow[d,"h_U"]\\
\mathbf P^1&U\arrow[l]
\end{tikzcd}\]
其中 $h=f|_{X'}$。以 $L$ 表示异于 $p$ 的所有\Term{prime-number}{素数}的\Term{set}{集合}，$\bar\eta$ 表示位于 $\mathbf P^1$ 的\Term{generic-point}{泛点} $\eta$ 上方的一个\Term{geometric-point}{几何点}。验证 SGA 1 XIII 4.8 的假设。因 $f$ 有\Term{section}{截面} $\sigma$，$h$ 也有。先证 SGA 1 XIII 4.7 的条件 a)：因 $h$ 的\Term{fiber}{纤维}\Term{geometrically-integral}{几何整}，$h$ 是 $0$-\Term{acyclic}{无环}且\Term{local-scheme}{局部} $0$-\Term{acyclic}{无环}的（SGA 4 XV 4.1、1.16）。还须说明，对 $X'$ 上任意\Term{locally-constant}{局部常值}\Term{set-sheaf}{集合层} $F$，$(F,h)$ 在\Term{dimension}{维数} $\leq0$ 上\Term{cohomologically-proper}{上同调固有}。若 $\xi$ 是 $\mathbf P^1$ 的\Term{closed-point}{闭点}，$Z$ 是 $\operatorname{Spec}\mathscr O_{\mathbf P^1,\xi}$ 在 $\bar\eta$ 中的\Term{integral-closure}{整闭包}，则\Term{scheme}{概形} $X'_Z=X'\times_{\mathbf P^1}Z$ 在 $X'$ 的\Term{closed-point}{闭点}处满足 $(S_2)$，在其余各点\Term{regular}{正则}，因此\Term{normal}{正规}；所以 $X'_Z$ 的一个\Term{etale-covering}{平展覆盖}，若在 $X'_{\bar\eta}$ 上的\Term{restriction}{限制}\Term{connected}{连通}，就必然\Term{connected}{连通}。这证明所需断言，故 a) 成立。由于 $X'_U$ 是 $S'_U$ 中一个相对于 $U$ \Term{normal-crossings}{正规交叉}的\Term{divisorial}{除子的}补集，$h_U$ 对 $L$ \Term{local-scheme}{局部} $1$-无球（SGA 4 XV 2.1），且对 $X'_U$ 上任意\Term{finite}{有限}\Term{constant}{常值} $L$-\Term{group-sheaf}{群层} $F$，$(F,h_U)$ 在\Term{dimension}{维数} $\leq1$ 上\Term{cohomologically-proper}{上同调固有}（SGA 1 XIII 2.4），故 b) 成立。
\SourcePage{39}
由于 $\mathbf P^1$ \Term{normal}{正规}，$T=\mathbf P^1-U$ \Term{codimension}{余维}为 $1$，故
\[\operatorname{prof\,\text{ét}}_T(S)\geq2,\]
从而 c) 成立。由\Term{purity-theorem}{纯性定理}，在 $X'$ 的每一\Term{closed-point}{闭点}都有
\[\operatorname{prof\,hop}_x(X')\geq3,\]
故 e) 成立。最后，对每个 $t\in\mathbf P^1$，$h$ 在 $\sigma(t)$ 处\Term{smooth}{光滑}，因而在此点对 $L$ \Term{local-scheme}{局部} $1$-无球，所以 d) 成立。

由 SGA 1 XIII 4.7，若 $a$ 是 $X'_{\bar\eta}$ 的\Term{geometric-point}{几何点}，则有\Term{exact-sequence}{正合列}
\[1\longrightarrow\Pi_1^{(p)}(X'_{\bar\eta},a)
\longrightarrow\Pi'_1(X'_U,a)
\longrightarrow\Pi_1(U,a)\longrightarrow1,\]
而\Term{section}{截面} $\sigma$ 的选择给出该\Term{exact-sequence}{正合列}的一个\Term{splitting}{分裂}。记 $K$ 为 $\Pi_1(U,a)$ 在 $\operatorname{Aut}(\Pi_1^{(p)}(X'_{\bar\eta},a))$ 中的像。$K$ \Term{action}{作用}下的\Term{coinvariants}{余不变量}群 $\Pi_1^{(p)}(X'_{\bar\eta},a)_K$ 是将 $\Pi_1^{(p)}(X'_{\bar\eta},a)$ 除以由形如 $x^{-1}qx$ 的元素生成的闭\Term{normal-subgroup}{正规子群}而得的商，其中 $x\in\Pi_1^{(p)}(X'_{\bar\eta},a)$、$q\in K$。由同上 (4.7.4)，有\Term{isomorphism}{同构}
\[\Pi_1^{(p)}(X',a)\simeq\Pi_1^{(p)}(X'_{\bar\eta},a)_K.\]
令 $\mathbf P^1-U=\coprod_{1\leq i\leq n}\{t_i\}$；对每个 $i$，令 $\overline T_i$ 是 $\mathbf P^1$ 在 $t_i$ 上方某个\Term{geometric-point}{几何点} $\bar t_i$ 处的\Term{strict-localization}{严格局部化}，$s_i$ 是 $\overline T_i$ 的\Term{generic-point}{泛点}。$U$ 的一个\Term{etale-covering}{平展覆盖}，若在每个 $s_i$ 上的\Term{restriction}{限制}平凡，就延拓为 $\mathbf P^1$ 的\Term{etale-covering}{平展覆盖}，因而本身平凡。对每个 $i$，设 $\overline{k(s_i)}$ 是 $k(s_i)$ 的\Term{algebraic-closure}{代数闭包}；上述断言即说，$\Pi_1(U,a)$ 等于由各 Galois 群 $\operatorname{Gal}(\overline{k(s_i)}/k(s_i))$ 在其中的像生成的闭\Term{normal-subgroup}{正规子群}。由第 I 讲 5.3，$\operatorname{Gal}(\overline{k(s_i)}/k(s_i))$ 在 $\operatorname{Aut}(\Pi_1^{(p)}(X'_{\bar\eta},a))$ 中的像 $K_i$ \Term{finite}{有限}。因为 $K$ 是包含所有 $K_i$ 的最小闭\Term{normal-subgroup}{正规子群}，且 $\Pi_1^{(p)}(X'_{\bar\eta},a)$ \Term{topologically-finitely-presented}{在拓扑意义下具有有限表示}（SGA 1 XIII 3.2），所以 $\Pi_1^{(p)}(X'_{\bar\eta},a)_K$ \Term{topologically-finitely-presented}{在拓扑意义下具有有限表示}，从而与之\Term{isomorphism}{同构}的 $\Pi_1^{(p)}(X)$ 亦然。
\end{proof}
\SourcePage{40}
\textbf{注记 2.3.2。} 若 $X$ 是 $k$ 上\Term{smooth}{光滑}\Term{projective-scheme}{射影概形}中一个\Term{normal-crossings-divisor}{正规交叉除子}的开补集，则证明\Term{theorem}{定理} 2.3.1 不需要\Term{resolution-singularities}{奇点消解}。
\subsection*{参考文献}
\begin{enumerate}
\item[{[1]}] S. Abhyankar，\emph{Local uniformization of algebraic surfaces over ground field of characteristic $p=0$}，Amer. J. of Math.，第 63 卷，1956。
\item[{[2]}] H. Hironaka，\emph{Resolution of singularities of an algebraic variety over a field of characteristic zero}，Ann. of Math.，第 79 卷，1964。
\item[{[3]}] M. Raynaud，\emph{Théorème de Lefschetz en cohomologie des faisceaux cohérents et en cohomologie étale}，学位论文，待刊。
\end{enumerate}

\SGAchapter{VI}{形式形变理论}{Formal deformation theory}
\SourcePage{41}
\begin{center}D. S. Rim\end{center}
\SGAsection{1}{形式存在定理}
回顾\Term{cofibered-category}{余纤维范畴}的概念（SGA 1 IV）。固定一个\Term{category}{范畴} $C$。按定义，$C$ 上的\Term{cofibered-category}{余纤维范畴}是一个\Term{category}{范畴} $\mathcal A$ 连同\Term{functor}{函子} $P:\mathcal A\to C$，满足：

\textbf{公理 1。} 对 $C$ 中任意\Term{morphism}{态射} $f$ 和 $\mathcal A$ 中任意\Term{object}{对象} $a$，若 $P(a)$ 为 $f$ 的源，则存在一个\Term{cocartesian}{余笛卡儿}的 $f$-\Term{morphism}{态射} $\alpha:a\to b$；即对任意 $f$-\Term{morphism}{态射} $\alpha':a\to b'$，存在唯一的 $T$-\Term{morphism}{态射} $\tau:b\to b'$ 使 $\alpha'=\tau\alpha$，其中 $T$ 为 $f$ 的靶。

\textbf{公理 2。} $\mathcal A$ 中\Term{cocartesian-morphism}{余笛卡儿态射}的复合仍为\Term{cocartesian-morphism}{余笛卡儿态射}。

从定义立即得出\Term{cofibered-category}{余纤维范畴} $P:\mathcal A\to C$ 的下列性质。

\textbf{消去律。} 设 $a\rightrightarrows a'$ 为两个\Term{cocartesian-morphism}{余笛卡儿态射} $\alpha,\beta$，且 $P(\alpha)=P(\beta)$。若存在\Term{cocartesian-morphism}{余笛卡儿态射} $\gamma:b\to a$ 使 $\alpha\gamma=\beta\gamma$，则 $\alpha=\beta$。

\textbf{\Term{resolution}{分解}。} 给定\Term{cocartesian-morphism}{余笛卡儿态射} $\alpha:a\to a'$、$\beta:a\to a''$，存在\Term{cocartesian-morphism}{余笛卡儿态射} $\gamma:a'\to a''$ 使 $\gamma\alpha=\beta$，当且仅当 $C$ 中存在 $f:P(a')\to P(a'')$ 使 $fP(\alpha)=P(\beta)$；此外 $\gamma$ 唯一。

设 $\mathcal A$ 是 $C$ 上的\Term{cofibered-category}{余纤维范畴}。对 $C$ 中每个\Term{object}{对象} $S$，以 $\mathcal A(S)$ 表示 $\mathcal A$ 的\Term{subcategory}{子范畴}，其\Term{object}{对象}为满足 $P(a)=S$ 的\Term{object}{对象} $a$，且
\[\operatorname{Hom}_{\mathcal A(S)}(a,a')=
\{\alpha\in\operatorname{Hom}_{\mathcal A}(a,a')\mid P(\alpha)=\operatorname{id}_S\}.\]
$C$ 上的\Term{cofibered-groupoid}{余纤维群胚}是这样一个\Term{cofibered-category}{余纤维范畴} $\mathcal A$：对 $C$ 的每个\Term{object}{对象} $S$，$\mathcal A(S)$ 都是\Term{groupoid}{群胚}。
\SourcePage{42}
回顾，\Term{groupoid}{群胚}是其中每个\Term{morphism}{态射}均为\Term{isomorphism}{同构}的\Term{category}{范畴}。因此，$C$ 上的\Term{cofibered-groupoid}{余纤维群胚}就是一个\Term{cofibered-category}{余纤维范畴} $P:\mathcal A\to C$，其中 $\mathcal A$ 的每个\Term{morphism}{态射}均为\Term{cocartesian-morphism}{余笛卡儿态射}。此外，$\alpha\in\operatorname{Fl}(\mathcal A)$ 是 $\mathcal A$ 中的\Term{isomorphism}{同构}，当且仅当 $P(\alpha)$ 是\Term{isomorphism}{同构}。

本节以下固定\Term{complete}{完备} \Term{noetherian-local-ring}{Noether 局部环} $\Lambda$，其\Term{residue-field}{剩余域}为 $k$，并令 $C_\Lambda$ 表示具有同一\Term{residue-field}{剩余域} $k$ 的 Artin \Term{local-scheme}{局部} $\Lambda$-\Term{algebra}{代数}\Term{category}{范畴}，\Term{morphism}{态射}为 $\Lambda$ 上的\Term{local-homomorphism}{局部同态}。从 1.2 起，对 $C_\Lambda$ 上的\Term{cofibered-groupoid}{余纤维群胚} $\mathcal A$ 总假定 $\mathcal A(k)=\{e\}$；这里 $\{e\}$ 表示对任意两个\Term{object}{对象} $a,b$，$\operatorname{Hom}_{\{e\}}(a,b)$ 恰含一个元素的\Term{category}{范畴}。若 $\mathcal A$ 中的\Term{morphism}{态射} $a'\xrightarrow\alpha a$ 满足 $P(a')=P(a)$ 且 $P(\alpha)=1_{P(a)}$，就简称其为 $C_\Lambda$-\Term{morphism}{态射}。若 $P(\alpha):P(a')\to P(a)$ 在 $C_\Lambda$ 中\Term{surjectivity}{满射}，则也略滥用术语称 $\alpha$ 为\Term{surjectivity}{满射}。

\textbf{例 1.1(a)。} 设 $F:C\to(\text{范畴})$ 为\Term{functor}{函子}。如下定义\Term{category}{范畴} $\mathcal F$：其\Term{object}{对象}为二元组 $(S,a)$，其中 $S\in\operatorname{ob}(C)$、$a\in\operatorname{ob}F(S)$；\Term{morphism}{态射} $(S,a)\to(S',a')$ 是二元组 $(f,u')$，其中 $f:S\to S'$ 是 $C$ 中的\Term{morphism}{态射}，$u':F(f)a\to a'$ 是 $F(S')$ 中的\Term{morphism}{态射}。两个\Term{morphism}{态射}
\[(S,a)\xrightarrow{(f,u')}(S',a')\xrightarrow{(g,u'')}(S'',a'')\]
的复合为 $(gf,u''\cdot(F(g)u')):(S,a)\to(S'',a'')$。对由 $(S,a)\mapsto S$ 定义的\Term{functor}{函子} $P:\mathcal F\to C$，$\mathcal F$ 成为 $C$ 上的\Term{cofibered-category}{余纤维范畴}。若每个 $F(S)$ 都是\Term{groupoid}{群胚}，则 $\mathcal F$ 是 $C$ 上的\Term{cofibered-groupoid}{余纤维群胚}。特别地，\Term{functor}{函子} $F:C\to(\text{群})$ 给出余\Term{fiber}{纤维}群，\Term{functor}{函子} $F:C\to(\mathrm{Ens})$ 给出\Term{discrete}{离散}\Term{cofibered-groupoid}{余纤维群胚}。还应注意，若 $F,G:C\to(\text{范畴})$ 为\Term{functor}{函子}，\Term{natural-transformation}{自然变换} $\varphi:F\to G$ 在 $C$ 上诱导\Term{cocartesian-functor}{余笛卡儿函子} $\Phi:\mathcal F\to\mathcal G$。

\textbf{例 1.1(b)。} 设 $X,Y$ 是 $\Lambda$ 上的\Term{scheme}{概形}，且固定\Term{isomorphism}{同构}
\SourcePage{43}
\[X\otimes_Ak\xrightarrow\tau Y\otimes_Ak.\]
\Term{functor}{函子} $\operatorname{Isom}_{X,Y}:C_\Lambda\to(\text{群胚})$，把 $S$ 映到满足 $\sigma\otimes_Ak=\tau$ 的\Term{isomorphism}{同构} $\sigma:X\otimes_AS\xrightarrow\sim Y\otimes_AS$ 的\Term{set}{集合}，给出\Term{cofibered-groupoid}{余纤维群胚} $P:\operatorname{Isom}_{X,Y}\to C_\Lambda$。

\textbf{例 1.1(c)。} 设 $X_0$ 是固定的 $k$-\Term{scheme}{概形}，以 $\mathcal M_{X_0}$ 表示由 $X_0$ 在 $C_\Lambda$ 上的\Term{deformation}{形变}构成的\Term{category}{范畴}（见第 4 节）。因此其\Term{object}{对象}是二元组 $(R,X)$，其中 $R\in C_\Lambda$，$X$ 是 $X_0$ 到 $R$ 的一个\Term{deformation}{形变}。由 $(R,X)\mapsto R$ 给出的\Term{functor}{函子} $P:\mathcal M_{X_0}\to C_\Lambda$ 定义\Term{cofibered-groupoid}{余纤维群胚}。

\textbf{例 1.1(d)。} 对任意\Term{groupoid}{群胚} $\mathcal X$，以 $[\mathcal X]$ 表示由其\Term{object}{对象}的\Term{isomorphism-class}{同构类}构成的\Term{discrete}{离散}\Term{groupoid}{群胚}。若 $P:\mathcal F\to C_\Lambda$ 为\Term{cofibered-groupoid}{余纤维群胚}，则得到\Term{functor}{函子} $[\mathcal F]:C_\Lambda\to(\mathrm{Ens})$，定义为 $[\mathcal F](S)=[\mathcal F(S)]$，进而得到\Term{discrete}{离散}\Term{cofibered-groupoid}{余纤维群胚}，也略滥用记号记作 $[\mathcal F]$。
\begin{SGAresult}{定义}{1.2}
称 $C_\Lambda$ 上的\Term{cofibered-groupoid}{余纤维群胚} $\mathcal A$ 为\Term{semi-homogeneous}{半齐性的}，若满足：
\begin{enumerate}
\item[(1)] $\mathcal A$ 中的每个实线图，其中 $a'\to a$ \Term{surjectivity}{满射}，都可补上虚线\Term{arrow}{箭头}成为\Term{commutative-diagram}{交换图}：
\[\begin{tikzcd}a''\arrow[r,dashed]\arrow[d,dashed]&a'\arrow[d]\\a_1\arrow[r]&a.\end{tikzcd}\]
\item[(2)] 设 $R\times_k k[\epsilon]\xrightarrow{\pi_1}R$、$R\times_k k[\epsilon]\xrightarrow{\pi_2}k[\epsilon]$ 为\Term{projection}{投影}，其中 $k[\epsilon]$ 是 $k$ 上的\Term{dual-numbers}{对偶数环}。设 $\alpha:a'\to a$、$\beta:b'\to a$ 为 $\pi_1$-\Term{morphism}{态射}。则存在 $C_\Lambda$-\Term{morphism}{态射} $\gamma:a'\to b'$ 使 $\alpha=\beta\gamma$，当且仅当 $(\pi_2)_*a'\simeq(\pi_2)_*b'$。
\end{enumerate}
\end{SGAresult}
\textbf{注记 1.3(a)。} 稍后在 2.6(b) 中将用 $2$-\Term{category}{范畴}的语言重新表述这个概念。注意，若 $\mathcal A$ \Term{semi-homogeneous}{半齐性}，则 $[\mathcal A]$ 也\Term{semi-homogeneous}{半齐性}，且 1.2(2) 特别蕴含\Term{canonical-map}{典范映射}
\SourcePage{44}
\[[\mathcal A(R\times_k k[\epsilon])]\longrightarrow[\mathcal A(R)]\times[\mathcal A(k[\epsilon])]\]
为\Term{bijection}{双射}。

\textbf{注记 1.3(b)。} 考虑 $C_\Lambda$ 上\Term{cofibered-category}{余纤维范畴} $\mathcal A$ 中的\Term{commutative-diagram}{交换图}
\[\begin{tikzcd}&a'\arrow[dl,"v_1"']\arrow[dr,"v_2"]&\\a_1\arrow[dr,"u_1"']&&a_2\arrow[dl,"u_2"]\\&a&\end{tikzcd}\]
令 $R=P(a)$、$R'=P(a')$、$R_i=P(a_i)$（$i=1,2$）。把 $a'$ 沿映射 $(P(v_1),P(v_2)):R'\to R_1\times_R R_2$ 的“\Term{direct-image}{正像}”记为 $a''$，就得到\Term{commutative-diagram}{交换图}
\[\begin{tikzcd}&a'\arrow[ddl,bend right,"v_1"']\arrow[ddr,bend left,"v_2"]\arrow[d]&\\&a''\arrow[dl]\arrow[dr]&\\a_1\arrow[dr,"u_1"']&&a_2\arrow[dl,"u_2"]\\&a&\end{tikzcd}\]
因此 1.2(1) 可换为：

(1)$'$：每个 $\mathcal A$ 中位于 $R_1\to R\leftarrow R_2$ 上方的图 $a_1\to a\leftarrow a_2$，其中 $R_2\to R$ \Term{surjectivity}{满射}，都能补为\Term{commutative-diagram}{交换图}
\[\begin{tikzcd}&a'\arrow[dl]\arrow[dr]&\\a_1\arrow[dr]&&a_2\arrow[dl]\\&a&\end{tikzcd}
\qquad\text{位于}\qquad
\begin{tikzcd}&R_1\times_RR_2\arrow[dl,"\pi_1"']\arrow[dr,"\pi_2"]&\\R_1\arrow[dr]&&R_2\arrow[dl]\\&R&\end{tikzcd}\]
之上，其中 $\pi_i$（$i=1,2$）为\Term{projection}{投影}。

\textbf{注记 1.3(c)。} 若 $\mathcal A$ 是 $C_\Lambda$ 上的\Term{semi-homogeneous}{半齐性}\Term{cofibered-groupoid}{余纤维群胚}，则 $[\mathcal A(k[\epsilon])]$ \Term{canonical}{典范}地具有 $k$-\Term{vector-space}{向量空间}结构，其中 $k[\epsilon]$ 是 $k$ 上的\Term{dual-numbers}{对偶数环}。事实上，令 $V$ 表示\Term{finite-dimensional}{有限维} $k$-\Term{vector-space}{向量空间}\Term{category}{范畴}。
\SourcePage{45}
任意保持积的\Term{functor}{函子} $H:V\to(\mathrm{Ens})$ 都\Term{canonical}{典范}地经由下图\Term{resolution}{分解}：
\[\begin{tikzcd}V\arrow[rr,"H"]\arrow[dr,"\widetilde H"']&&(\mathrm{Ens})\\&V\arrow[ur]&\end{tikzcd}\]
其中 $V\to(\mathrm{Ens})$ 为\Term{forgetful-functor}{遗忘函子}。另一方面，由 $W\mapsto k[W]=k\oplus W$、$W^2=0$ 给出的\Term{functor}{函子} $G:V\to C_\Lambda$ 把 $V$ 中的积变为 $C_\Lambda$ 中在 $k$ 上的\Term{fiber-product}{纤维积}。若 $\mathcal A$ 是\Term{semi-homogeneous}{半齐性}\Term{cofibered-groupoid}{余纤维群胚}，则由 1.3(a)，
\[[\mathcal A(k[V]\times_k k[W])]\longrightarrow[\mathcal A(k[V])]\times[\mathcal A(k[W])]\]
是\Term{bijection}{双射}，故\Term{functor}{函子} $V\mapsto[\mathcal A(k[V])]$ 与积交换，结论得证。
\begin{SGAresult}{定义}{1.4}
(1) 设 $R\in\operatorname{ob}(C_\Lambda)$。$C_\Lambda$ 中 $R$ 的一个\Term{small-extension}{小扩张}是\Term{surjectivity}{满射} $f:R'\to R$，满足 $(\operatorname{Ker}f)^2=0$ 且 $\operatorname{Ker}f\simeq k$。换言之，$f:R'\to R$ 是\Term{small-extension}{小扩张}，当且仅当对某个 $t\in R'$、$t^2=0$，序列
\[0\longrightarrow k\xrightarrow{t}R'\xrightarrow fR\longrightarrow0\]
\Term{exact}{正合}。

(2) 设 $\mathcal A$ 是 $C_\Lambda$ 上的\Term{cofibered-groupoid}{余纤维群胚}，$a$ 是其\Term{object}{对象}。$a$ 的一个\Term{small-prolongation}{小延拓}是\Term{arrow}{箭头} $\alpha:a'\to a$，使 $P(\alpha)$ 是 $C_\Lambda$ 中的\Term{small-extension}{小扩张}。若对 $\mathcal A$ 中 $a$ 的每个\Term{small-prolongation}{小延拓} $a'_1\to a$，都有\Term{commutative-diagram}{交换图}
\[\begin{tikzcd}a'\arrow[rr]\arrow[dr,"\alpha"']&&a'_1\arrow[dl]\\&a&\end{tikzcd}\]
则称 $\alpha:a'\to a$ “对 $a$ 的\Term{small-prolongation}{小延拓}为通用的”。
\end{SGAresult}
\begin{SGAresult}{命题}{1.5}
设 $\mathcal A$ 是 $C_\Lambda$ 上的\Term{semi-homogeneous}{半齐性}\Term{cofibered-groupoid}{余纤维群胚}。

(1) 设 $R\times_k k[\epsilon]\xrightarrow{\pi_1}R$、$R\times_k k[\epsilon]\xrightarrow{\pi_2}k[\epsilon]$ 为 $C_\Lambda$ 中的\Term{projection}{投影}，$\alpha:a'\to a$ 为 $\pi_1$-\Term{morphism}{态射}。
\SourcePage{46}
则存在\Term{morphism}{态射} $\beta:a\to a'$ 使 $\alpha\beta=1_a$，当且仅当存在\Term{morphism}{态射} $a\to(\pi_2)_*a'$。

(2) 设 $f:R'\to R$ 是 $C_\Lambda$ 中的\Term{small-extension}{小扩张}，并考虑\Term{commutative-diagram}{交换图}
\[\begin{tikzcd}&a''\arrow[dl,"\alpha'"']\arrow[dr,"\alpha'_1"]&\\a'\arrow[dr,"\alpha"']&&a'_1\arrow[dl,"\alpha_1"]\\&a&\end{tikzcd}
\quad\text{位于}\quad
\begin{tikzcd}&R'\times_RR'\arrow[dl,"\pi_1"']\arrow[dr,"\pi_2"]&\\R'\arrow[dr,"f"']&&R'\arrow[dl,"f"]\\&R&\end{tikzcd}\]
之上。定义\Term{canonical-map}{典范映射}
\[\bar f:R'\times_RR'\longrightarrow k[\operatorname{Ker}f],\qquad
\bar f(r_1,r_2)=r_1(0)+(r_1-r_2).\]
若存在\Term{morphism}{态射} $a'\to(\bar f)_*(a'')$，则存在 $\gamma:a'\to a'_1$，使 $\alpha=\alpha_1\gamma$。
\end{SGAresult}
\begin{proof}
(1) “仅当”显然。现设有\Term{morphism}{态射} $\varphi:a\to(\pi_2)_*a'$。令 $h:R\to R\times_k k[\epsilon]$ 为 $h=1\times P(\varphi)$。复合 $R\xrightarrow hR\times_k k[\epsilon]\xrightarrow{\pi_1}R$ 为恒等，故有\Term{morphism}{态射}
\[a\xrightarrow\gamma h_*a\xrightarrow{\bar\gamma}a,\qquad
\bar\gamma\gamma=1_a,\quad P(\gamma)=h,\ P(\bar\gamma)=\pi_1.\]
又 $\pi_2h=P(\varphi)$，故有 $\pi_2$-\Term{morphism}{态射} $h_*a\to(\pi_2)_*a'$。这说明 $(\pi_2)_*(h_*a)=(\pi_2)_*a'$；因此由 1.2(2)，有 $C_\Lambda$-\Term{morphism}{态射} $u:h_*a\to a'$，使 $\alpha u=\bar\gamma$。令 $\beta=u\gamma$，便有 $\alpha\beta=\alpha u\gamma=\bar\gamma\gamma=1_a$。

(2) 映射 $1\times\bar f:R'\times_RR'\to R'\times_k k[\operatorname{Ker}f]$ 是与第一因子\Term{projection}{投影}相容的\Term{isomorphism}{同构}，并且 $\pi_2(1\times\bar f)=\bar f$，其中 $\pi_2:R'\times_k k[\operatorname{Ker}f]\to k[\operatorname{Ker}f]$ 为\Term{projection}{投影}。因 $R'\xrightarrow fR$ 是\Term{small-extension}{小扩张}，$k[\operatorname{Ker}f]\simeq k[\epsilon]$。因此，若存在\Term{morphism}{态射} $a'\to(\bar f)_*a''=(\pi_2)_*[(1\times\bar f)_*a'']$，则由 (1)，存在 $\gamma':a'\to a''$，使 $\alpha'\gamma'=1_{a'}$。令 $\gamma=\alpha'_1\gamma'$，就有 $\alpha_1\gamma=\alpha$。
\end{proof}
\SourcePage{47}
\begin{SGAresult}{推论}{1.6}
设 $\mathcal A$ 是 $C_\Lambda$ 上的\Term{semi-homogeneous}{半齐性}\Term{cofibered-groupoid}{余纤维群胚}，并考虑其中的图
\[\begin{tikzcd}a'\arrow[dr,"\alpha"']&&a_1\arrow[dl,"\alpha_1"]\\&a&\end{tikzcd}\]
其中 $\alpha_1$ 是\Term{small-prolongation}{小延拓}。假设 $a$ 对 $e$ 的\Term{small-prolongation}{小延拓}为通用的。则 $\mathcal A$ 中存在\Term{morphism}{态射} $\rho:a'\to a_1$ 使 $\alpha=\alpha_1\rho$，当且仅当 $C_\Lambda$ 中存在\Term{morphism}{态射} $f:P(a')\to P(a_1)$ 使 $P(\alpha)=P(\alpha_1)f$。
\end{SGAresult}
\begin{proof}
设 $P(\alpha)=P(\alpha_1)f$，其中 $f:P(a')\to P(a_1)$。将 $a'$ 换成沿 $f$ 的\Term{direct-image}{正像}，可且将假定 $P(\alpha)=P(\alpha_1)$。令 $R'=P(a_1)$、$R=P(a)$。由于 $\mathcal A$ \Term{semi-homogeneous}{半齐性}，得到\Term{commutative-diagram}{交换图}
\[\begin{tikzcd}&a''\arrow[dl,"\beta"']\arrow[dr,"\beta_1"]&\\a'\arrow[dr,"\alpha"']&&a_1\arrow[dl,"\alpha_1"]\\&a&\end{tikzcd}
\quad\text{位于}\quad
\begin{tikzcd}&R'\times_RR'\arrow[dl]\arrow[dr]&\\R'\arrow[dr]&&R'\arrow[dl]\\&R&\end{tikzcd}\]
之上。因 $a$ 对 $e$ 的\Term{small-prolongation}{小延拓}为通用的，$a'$ 亦如此；因此由 1.5(2) 即得结论。
\end{proof}
\begin{SGAresult}{定义}{1.7}
对 $C_\Lambda$ 中的每个\Term{object}{对象} $R$，令
\[\overline\Omega_R=M/(\mathfrak m_\Lambda R+M^2),\]
其中 $\mathfrak m_\Lambda,M$ 分别是 $\Lambda,R$ 的\Term{maximal-ideal}{极大理想}。这是\Term{finite-dimensional}{有限维} $k$-\Term{vector-space}{向量空间}，$R\mapsto\overline\Omega_R$ 定义\Term{functor}{函子} $\overline\Omega:C_\Lambda\to V$，其中 $V$ 为\Term{finite-dimensional}{有限维} $k$-\Term{vector-space}{向量空间}\Term{category}{范畴}。若 $\mathcal A$ 是 $C_\Lambda$ 上的\Term{cofibered-category}{余纤维范畴}，则\Term{composite-functor}{复合函子} $\mathcal A\xrightarrow P C_\Lambda\xrightarrow{\overline\Omega}V$ 也略滥用记号记为 $\overline\Omega$。
\end{SGAresult}
\begin{SGAresult}{命题}{1.8}
设 $\mathcal A$ 是 $C_\Lambda$ 上的\Term{semi-homogeneous}{半齐性}\Term{cofibered-groupoid}{余纤维群胚}。给定 $\mathcal A$ 中一个对 $e$ 的\Term{small-prolongation}{小延拓}为通用的\Term{object}{对象} $a$，则存在 $a'\to a$，使：
\SourcePage{48}
\begin{enumerate}
\item[(i)] $a'$ 对 $a$ 的\Term{small-prolongation}{小延拓}为通用的；
\item[(ii)] $P(a')\to P(a)$ \Term{surjectivity}{满射}，且其核被 $P(a)$ 的\Term{maximal-ideal}{极大理想}\Term{annihilate}{零化}；
\item[(iii)] $\overline\Omega_{a'}\to\overline\Omega_a$ 为\Term{isomorphism}{同构}。
\end{enumerate}
\end{SGAresult}
\begin{proof}
令 $R=P(a)$，并取\Term{exact-sequence}{正合列} $0\to J\to S\to R\to0$，其中 $S$ 是 $\Lambda$ 上\Term{finite}{有限}个变量的\Term{formal-power-series-ring}{形式幂级数环}，且 $J\subset M^2$，$M$ 为 $S$ 的\Term{maximal-ideal}{极大理想}。以 $\Sigma$ 表示满足下列条件的 $S$ 的\Term{ideal}{理想} $J_1$ 的\Term{set}{集合}：(i) $J\supset J_1\supset MJ$；(ii) $\mathcal A$ 中存在 $\pi_1$-\Term{morphism}{态射} $a_1\to a$，其中 $\pi_1:S/J_1\to R$ 是\Term{canonical-map}{典范映射}。$\mathcal A$ 的性质 1.2(1) 加上 $J/MJ$ 是\Term{finite-dimensional}{有限维} $k$-\Term{vector-space}{向量空间}，蕴含 $\Sigma$ 中存在最小\Term{ideal}{理想}。记其为 $J'$，并选取 $\pi'$-\Term{morphism}{态射} $a'\xrightarrow\alpha a$，其中 $\pi':S/J'\to R$ 是\Term{canonical-map}{典范映射}。因 $J'\subset J\subset M^2$，$\overline\Omega_{a'}\to\overline\Omega_a$ 是\Term{isomorphism}{同构}。断言 $a'\xrightarrow\alpha a$ 对 $a$ 的\Term{small-prolongation}{小延拓}为通用的。事实上，设 $a_1\to a$ 是一个\Term{small-prolongation}{小延拓}。令 $R_1=P(a_1)$，得到行\Term{exact}{正合}的\Term{commutative-diagram}{交换图}
\[\begin{tikzcd}
0\arrow[r]&J/MJ\arrow[r]\arrow[d,"h"]&S/MJ\arrow[r]\arrow[d,"h_1"]&R\arrow[r]\arrow[d,equal]&0\\
0\arrow[r]&k\arrow[r]&R_1\arrow[r,"\pi_1"']&R\arrow[r]&0.
\end{tikzcd}\]
若 $h=0$，则 $\pi_1:R_1\to R$ 有\Term{section}{截面}，因而 $h_1$ 诱导 $f:S/J'\to R_1$。若 $h\ne0$，则 $h_1:S/MJ\to R_1$ \Term{surjectivity}{满射}，所以由\Term{ideal}{理想} $J'$ 的定义，$h_1$ 诱导 $f:S/J'\to R_1$。因此两种情形下均有\Term{commutative-diagram}{交换图}
\[\begin{tikzcd}S/J'\arrow[rr,"f"]\arrow[dr,"\pi'"']&&R_1\arrow[dl,"\pi_1"]\\&R&\end{tikzcd}\]
因 $a$ 对 $e$ 的\Term{small-prolongation}{小延拓}为通用的，由 1.6 得到\Term{commutative-diagram}{交换图}
\SourcePage{49}
\[\begin{tikzcd}a'\arrow[rr,"\rho"]\arrow[dr]&&a_1\arrow[dl]\\&a&\end{tikzcd}\]
\end{proof}
设 $P:\mathcal A\to C_\Lambda$ 是\Term{cofibered-category}{余纤维范畴}（其\Term{fiber}{纤维}为\Term{groupoid}{群胚}）。显然得到\Term{cofibered-category}{余纤维范畴}（其\Term{fiber}{纤维}为\Term{groupoid}{群胚}）
\[\operatorname{pro}(P):\operatorname{pro}(\mathcal A)\longrightarrow\operatorname{pro}(C_\Lambda)\]
（关于投\Term{object}{对象}\Term{category}{范畴} $\operatorname{pro}(\mathcal A)$ 的定义，见 A. Grothendieck，Bourbaki 讨论班第 195 讲），并有\Term{commutative-diagram}{交换图}
\[\begin{tikzcd}\mathcal A\arrow[r]\arrow[d,"P"']&\operatorname{pro}(\mathcal A)\arrow[d,"\operatorname{pro}(P)"]\\C_\Lambda\arrow[r]&\operatorname{pro}(C_\Lambda).\end{tikzcd}\]
现令 $\widehat C_\Lambda$ 为 $\operatorname{pro}(C_\Lambda)$ 的\Term{full-subcategory}{全子范畴}，其\Term{object}{对象}为
\[\begin{split}
\operatorname{ob}(\widehat C_\Lambda)=\{(R_i)_{i\in\mathbf N}\mid\;&R_i\in\operatorname{ob}(C_\Lambda),\quad
\varphi_i:R_i\to R_{i-1}\text{ 对所有 }i\text{ 满射，}\\
&\text{且对充分大的 }i\text{，诱导同构 }
\mathfrak m_i/(\mathfrak m_\Lambda+\mathfrak m_i^2)
\xrightarrow\sim\mathfrak m_{i-1}/(\mathfrak m_\Lambda+\mathfrak m_{i-1}^2)\},
\end{split}\]
其中 $\mathbf N$ 为自然数的\Term{ordered-category}{有序范畴}，$\mathfrak m_i$ 是 $R_i$ 的\Term{maximal-ideal}{极大理想}。令 $\widehat{\mathcal A}$ 为 $\operatorname{pro}(\mathcal A)$ 的\Term{full-subcategory}{全子范畴}，其\Term{object}{对象}为
\[\operatorname{ob}(\widehat{\mathcal A})=
\{a\in\operatorname{ob}(\operatorname{pro}(\mathcal A))\mid
\operatorname{pro}(P)(a)\in\operatorname{ob}(\widehat C_\Lambda)\}.\]
于是得到\Term{cofibered-category}{余纤维范畴}（其\Term{fiber}{纤维}为\Term{groupoid}{群胚}）$\widehat P:\widehat{\mathcal A}\to\widehat C_\Lambda$，及\Term{commutative-diagram}{交换图}
\[\begin{tikzcd}\mathcal A\arrow[r,hook]\arrow[d,"P"']&\widehat{\mathcal A}\arrow[d,"\widehat P"]\\C_\Lambda\arrow[r,hook]&\widehat C_\Lambda.\end{tikzcd}\]
注意，$\widehat C_\Lambda$ \Term{canonical-equivalence}{典范等价}于具有固定\Term{residue-field}{剩余域} $k$、在 $\Lambda$ 上\Term{form}{形式}\Term{finite-type}{有限型}的\Term{complete}{完备}\Term{local-scheme}{局部} $\Lambda$-\Term{algebra}{代数}\Term{category}{范畴}。以下作此等同。
\begin{SGAresult}{定义}{1.9}
设 $\mathcal A$ 为 $C_\Lambda$ 上的\Term{cofibered-groupoid}{余纤维群胚}。若 $\widehat{\mathcal A}$ 中的\Term{object}{对象} $\xi$ 满足：对其中每个\Term{object}{对象} $a$ 均存在\Term{morphism}{态射} $\xi\to a$，则称 $\xi$ 为 $\widehat{\mathcal A}$ 的\Term{quasi-initial-object}{拟始对象}。若\Term{quasi-initial-object}{拟始对象} $\xi$ 对应的\Term{morphism}{态射} $\xi:h_R\to[\widehat{\mathcal A}]$ 诱导\Term{bijection}{双射}
\SourcePage{50}
\[\xi(k[\epsilon]):\operatorname{Hom}_{\Lambda\text{-alg}}(R,k[\epsilon])\longrightarrow[\mathcal A(k[\epsilon])],
\qquad R=P(\xi),\]
则称其为极小的。

若 $\widehat{\mathcal A}$ 中的\Term{object}{对象} $\xi$ 满足：$\mathcal A$ 中每个\Term{surjectivity}{满射} $\alpha:a'\to a$ 都诱导\Term{surjectivity}{满射}
\[\operatorname{Hom}_{\widehat{\mathcal A}}(\xi,a')\longrightarrow
\operatorname{Hom}_{\widehat{\mathcal A}}(\xi,a),\]
则称 $\xi$ 为 $\mathcal A$ 的\Term{versal-formal-object}{通用形式对象}。若\Term{versal-formal-object}{通用形式对象} $\xi$ 对应的\Term{morphism}{态射} $\xi:h_R\to[\widehat{\mathcal A}]$ 诱导\Term{bijection}{双射}
\[\xi(k[\epsilon]):\operatorname{Hom}_{\Lambda\text{-alg}}(R,k[\epsilon])\longrightarrow[\mathcal A(k[\epsilon])],
\qquad R=P(\xi),\]
则称其为极小的。
\end{SGAresult}
\begin{SGAresult}{命题}{1.10}
设 $\mathcal A$ 是 $C_\Lambda$ 上的\Term{cofibered-groupoid}{余纤维群胚}。
\begin{enumerate}
\item[(i)] $\mathcal A$ 的每个（极小）\Term{versal-formal-object}{通用形式对象}都是 $\widehat{\mathcal A}$ 的（极小）\Term{quasi-initial-object}{拟始对象}。
\item[(ii)] 若 $\widehat{\mathcal A}$ 中的\Term{object}{对象} $\xi$ 满足
\[\widetilde\xi(k[\epsilon]):\operatorname{Hom}_{\Lambda\text{-alg}}(R,k[\epsilon])\to[\mathcal A(k[\epsilon])]
\quad\text{为单射，其中 }P(\xi)=R,\]
则 $\xi$ 在 $\widehat{\mathcal A}$ 中的每个\Term{endomorphism}{自态射}都是\Term{automorphism}{自同构}。特别地，极小\Term{quasi-initial-object}{拟始对象}若存在，则在（非\Term{canonical}{典范}）\Term{isomorphism}{同构}意义下唯一。
\end{enumerate}
\end{SGAresult}
\begin{proof}
(i) 设 $\xi$ 是 $\mathcal A$ 的\Term{versal-formal-object}{通用形式对象}。须说明对 $\widehat{\mathcal A}$ 中任意\Term{object}{对象} $a$ 均有\Term{morphism}{态射} $\xi\to a$。由 $\widehat{\mathcal A}$ 的定义，存在 $\mathcal A$ 中的\Term{surjectivity}{满射}序列
\[\cdots\longrightarrow a_n\xrightarrow{\alpha_n}a_{n-1}\longrightarrow\cdots
\longrightarrow a_1\xrightarrow{\alpha_1}a_0,\qquad a=\varprojlim a_n.\]
假设已存在\Term{morphism}{态射} $\varphi_i:\xi\to a_i$，使 $i<n$ 时 $\alpha_i\varphi_i=\varphi_{i-1}$。因 $\xi$ 为\Term{versal-formal-object}{通用形式对象}，且 $\alpha_n:a_n\to a_{n-1}$ 在 $\mathcal A$ 中\Term{surjectivity}{满射}，得到 $\varphi_n:\xi\to a_n$ 使 $\alpha_n\varphi_n=\varphi_{n-1}$。令 $\varphi=\varprojlim\varphi_n$，便得 $\varphi:\xi\to a$。

(ii) 设 $\alpha:\xi\to\xi$ 是 $\widehat{\mathcal A}$ 中的\Term{morphism}{态射}。对每个 $\lambda\in\operatorname{Hom}_{\Lambda\text{-alg}}(R,k[\epsilon])$，有
\[\lambda_*\xi=\lambda_*\alpha_*\xi\simeq(\lambda\cdot P(\alpha))_*\xi.\]
由对 $\xi$ 的假设，对所有这样的 $\lambda$ 均有 $\lambda=\lambda\cdot P(\alpha)$，所以 $P(\alpha):R\to R$ \Term{surjectivity}{满射}。Noether 环的每个满\Term{endomorphism}{自同态}显然均为\Term{automorphism}{自同构}，因此 $P(\alpha)$ 是 $R$ 的\Term{automorphism}{自同构}，$\alpha$ 从而也是 $\xi$ 的\Term{automorphism}{自同构}。
\end{proof}
\begin{SGAresult}{\Term{theorem}{定理}}{1.11（Schlessinger）}
设 $\mathcal A$ 是 $C_\Lambda$ 上的\Term{cofibered-groupoid}{余纤维群胚}。
\SourcePage{51}
则 $\mathcal A$ 存在\Term{minimally-versal-formal-object}{极小通用形式对象}，当且仅当：
\begin{enumerate}
\item[(i)] $\mathcal A$ \Term{semi-homogeneous}{半齐性}；
\item[(ii)] $[\mathcal A(k[\epsilon])]$ 是\Term{finite-dimensional}{有限维} $k$-\Term{vector-space}{向量空间}（参见 1.3(c)）。
\end{enumerate}
\end{SGAresult}
\begin{proof}
\textbf{必要性。} 设 $\xi\in\widehat{\mathcal A}$ 为 $\mathcal A$ 的\Term{minimally-versal-formal-object}{极小通用形式对象}，并考虑 $\mathcal A$ 中的图
\[\begin{tikzcd}a_1\arrow[dr]&&a_2\arrow[dl,"\alpha"]\\&a&\end{tikzcd}\]
其中 $\alpha$ \Term{surjectivity}{满射}。因 $\xi$ 是 $\widehat{\mathcal A}$ 的\Term{quasi-initial-object}{拟始对象}，存在\Term{morphism}{态射} $\xi\to a_1$。由于 $\alpha$ \Term{surjectivity}{满射}，\Term{composition-morphism}{复合态射} $\xi\to a_1\to a$ \Term{liftable}{可提升}到 $a_2$，即得到\Term{commutative-diagram}{交换图}
\[\begin{tikzcd}&\xi\arrow[dl]\arrow[dr]&\\a_1\arrow[dr]&&a_2\arrow[dl,"\alpha"]\\&a&\end{tikzcd}\]
这证明条件 1.2(1)。现证明条件 1.2(2)。设 $S\times_k k[\epsilon]\xrightarrow{\pi_1}S$、$S\times_k k[\epsilon]\xrightarrow{\pi_2}k[\epsilon]$ 为\Term{projection}{投影}，$\alpha:a'\to a$、$\beta:b'\to a$ 为 $\mathcal A$ 中的 $\pi_1$-\Term{morphism}{态射}。由 $\xi$ 的定义，得到\Term{commutative-diagram}{交换图}
\[\begin{tikzcd}&\xi\arrow[dl,"u"']\arrow[dr,"v"]&\\a'\arrow[dr,"\alpha"']&&b'\arrow[dl,"\beta"]\\&a&\end{tikzcd}\]
特别地，$\pi_1P(u)=\pi_1P(v)$。现假设 $(\pi_2)_*a'\simeq(\pi_2)_*b'$。由 $\xi$ 的极小性，$\pi_2P(u)=\pi_2P(v)$，
\SourcePage{52}
所以
\[P(u)=\pi_1P(u)\times\pi_2P(u)=\pi_1P(v)\times\pi_2P(v)=P(v).\]
因 $\mathcal A$ 为\Term{cofibered-groupoid}{余纤维群胚}，得到 $\mathcal A(S\times_k k[\epsilon])$ 中唯一\Term{isomorphism}{同构} $\sigma:a'\to b'$，使 $\sigma u=v$。于是 $\beta\sigma u=\alpha u$；由于 $\mathcal A$ 是\Term{cofibered-category}{余纤维范畴}，故 $\beta\sigma=\alpha$。

\textbf{充分性。} 取 $\mathcal A(k[\epsilon])$ 中的\Term{object}{对象} $c_1,\ldots,c_n$，使其\Term{isomorphism-class}{同构类}构成 $[\mathcal A(k[\epsilon])]$ 的一个 $k$-基，并在 $\mathcal A(k[\epsilon]\times_k\cdots\times_k k[\epsilon])$ 中选取 $a_1$，使 $(\pi_j)_*a_1\simeq c_j$。则
\[\widetilde a(k[\epsilon]):\operatorname{Hom}_{\Lambda\text{-alg}}(R_1,k[\epsilon])
\longrightarrow[\mathcal A(k[\epsilon])],\qquad R_1=P(a_1),\]
为\Term{bijection}{双射}。特别地，$a_1\to e$ 对 $e$ 的\Term{small-prolongation}{小延拓}为通用的。选取对 $a_1$ 的\Term{small-prolongation}{小延拓}为通用的 $a_2\to a_1$，使 $\overline\Omega_{a_2}\to\overline\Omega_{a_1}$ 为\Term{bijection}{双射}。再取对 $a_2$ 的\Term{small-prolongation}{小延拓}为通用的 $a_3\to a_2$，使 $\overline\Omega_{a_3}\to\overline\Omega_{a_2}$ 为\Term{bijection}{双射}，依此继续，令 $\xi=\varprojlim a_n$。断言 $\xi$ 是 $\mathcal A$ 的\Term{minimally-versal-formal-object}{极小通用形式对象}。事实上，因所有 $n$ 对应的 $\overline\Omega_{a_n}\to\overline\Omega_{a_{n-1}}$ 均为\Term{bijection}{双射}，$\overline\Omega_\xi\to\overline\Omega_{a_1}$ 也是\Term{bijection}{双射}，故
\[\xi(k[\epsilon]):\operatorname{Hom}_{\Lambda\text{-alg}}(R,k[\epsilon])\to[\mathcal A(k[\epsilon])]\]
为\Term{bijection}{双射}。现在考虑图
\[\begin{tikzcd}&a'\arrow[d,"\alpha"]\\\xi\arrow[r,"h"']&a\end{tikzcd}\]
其中 $\alpha$ 在 $\mathcal A$ 中\Term{surjectivity}{满射}。须说明 $h$ \Term{liftable}{可提升}到 $a'$。为此，由归纳可假设 $a'$ 是 $a$ 的\Term{small-prolongation}{小延拓}。由于 $P(a)$ 是 Artin 环，对某个 $q$，$h:\xi\to a$ \Term{decomposable}{可分解}为 $\xi\to a_q\to a$。$\mathcal A$ \Term{semi-homogeneous}{半齐性}，故得到图
\[\begin{tikzcd}&a''\arrow[r]\arrow[d,"\alpha'"']&a'\arrow[d,"\alpha"]\\\xi\arrow[r]&a_q\arrow[r]&a\end{tikzcd}\]
其中 $\alpha'$ 也是\Term{small-prolongation}{小延拓}。于是有\Term{commutative-diagram}{交换图}
\[\begin{tikzcd}&&a''\arrow[d,"\alpha'"]\\\xi\arrow[r]&a_{q+1}\arrow[ur]\arrow[r]&a_q\end{tikzcd}\]
证明完毕。
\end{proof}
\textbf{注记 1.12。} 设
\SourcePage{53}
\[\cdots\longrightarrow a_n\longrightarrow a_{n-1}\longrightarrow\cdots
\longrightarrow a_2\longrightarrow a_1\longrightarrow a_0=e\]
是 $\mathcal A$ 中的无穷\Term{arrow}{箭头}序列，满足
\[\begin{cases}
a_i\text{ 对 }a_{i-1}\text{ 的小延拓为通用的},\\
P(a_i)\longrightarrow P(a_{i-1})\text{ 满射},\\
\overline\Omega_{a_i}\longrightarrow\overline\Omega_{a_{i-1}}\text{ 在 }i\gg0\text{ 时为双射}.
\end{cases}\]
则上述证明表明，$\varprojlim a_i$ 是 $\mathcal A$ 的\Term{versal-formal-object}{通用形式对象}。

\textbf{注记 1.13。} 设 $\mathcal A$ 是 $C_\Lambda$ 上的\Term{cofibered-groupoid}{余纤维群胚}。若 $\mathcal A$ 具有性质 1.11(i)、(ii)，则 $[\mathcal A]$ 也有这些性质；并且由 1.10(ii)，$\xi\in\operatorname{ob}(\widehat{\mathcal A})$ 是 $\mathcal A$ 的\Term{minimally-versal-formal-object}{极小通用形式对象}，当且仅当它是 $[\mathcal A]$ 的\Term{minimally-versal-formal-object}{极小通用形式对象}。

\textbf{注记 1.14。} 设 $\mathcal A$ 是 $C_\Lambda$ 上的\Term{cofibered-groupoid}{余纤维群胚}，且存在\Term{minimally-versal-formal-object}{极小通用形式对象} $\xi$。则 $R=P(\xi)$（它是 $\Lambda$ 上\Term{form}{形式}\Term{finite-type}{有限型}的\Term{complete}{完备}\Term{local-scheme}{局部} $\Lambda$-\Term{algebra}{代数}）的\Term{invariants}{不变量}自动成为 $C_\Lambda$ 上 $\mathcal A$ 的\Term{invariants}{不变量}。例如，若 $R$ 在 $\Lambda$ 上（\Term{form}{形式}）\Term{smooth}{光滑}，即为 $\Lambda$ 上的\Term{formal-power-series-algebra}{形式幂级数代数}，就称 $\mathcal A$ 在 $C_\Lambda$ 上（\Term{form}{形式}）\Term{smooth}{光滑}；若 $\operatorname{Spec}(R)$ \Term{irreducible}{不可约}，就称 $\mathcal A$ 在 $C_\Lambda$ 上\Term{irreducible}{不可约}，等等。下一节将引入 $C_\Lambda$ 上\Term{functor}{函子} $\mathcal A\to\mathcal B$ 的“\Term{smoothness}{光滑性}”，其中 $\mathcal A,\mathcal B$ 是 $C_\Lambda$ 上\Term{fiber}{纤维}为\Term{groupoid}{群胚}的\Term{cofibered-category}{余纤维范畴}。届时可知（参见 2.4），“$\mathcal A$ 在 $C_\Lambda$ 上（\Term{form}{形式}）\Term{smooth}{光滑}”\Term{equivalence}{等价}于“$P:\mathcal A\to C_\Lambda$ 是\Term{smooth-functor}{光滑函子}”，因而不会混淆。还定义 $\mathcal A$ 在 $C_\Lambda$ 上的\Term{dimension}{维数}为
\[\dim_{C_\Lambda}\mathcal A=\dim(k\otimes_\Lambda R).\]
为此注意下述事实：令 $C_k$ 为具有同一\Term{residue-field}{剩余域} $k$ 的 Artin \Term{local-scheme}{局部} $k$-\Term{algebra}{代数}\Term{category}{范畴}，则 $C_k\hookrightarrow C_\Lambda$ 为\Term{fully-faithful}{全忠实}\Term{embedding}{嵌入}，其（左）\Term{adjoint-functors}{伴随函子} $C_\Lambda\to C_k$ 为 $R\mapsto k\otimes_\Lambda R$。设 $\mathcal A|C_k$ 是 $P:\mathcal A\to C_\Lambda$ 沿 $C_k\hookrightarrow C_\Lambda$ 的拉回。显然，若 $\xi$ 是 $C_\Lambda$ 上 $\mathcal A$ 的\Term{minimally-versal-formal-object}{极小通用形式对象}，则 $k\otimes_\Lambda\xi$ 是 $C_k$ 上 $\mathcal A|C_k$ 的\Term{minimally-versal-formal-object}{极小通用形式对象}；其中 $\xi\to k\otimes_\Lambda\xi$ 为位于 $R\to k\otimes_\Lambda R$ 之上的\Term{cocartesian}{余笛卡儿}\Term{arrow}{箭头}。因此
\[\dim_{C_\Lambda}\mathcal A=\dim_{C_k}(\mathcal A|C_k).\]
\SourcePage{54}
\textbf{1.15。} 本节余下部分简要说明如何修改上述理论，以允许\Term{residue-field}{剩余域}\Term{group-extension}{扩张}。

设 $\Lambda,K$ 为\Term{complete}{完备} \Term{noetherian-local-ring}{Noether 局部环}，$K$ 为 $\Lambda$-\Term{algebra}{代数}。假设：
\begin{enumerate}
\item[(a)] 从 $\Lambda$ 到 $K$ 的自然映射为\Term{local-homomorphism}{局部同态}；
\item[(b)] $K$ 的\Term{residue-field}{剩余域} $K'$ 是 $\Lambda$ 的\Term{residue-field}{剩余域} $k$ 的\Term{finite-type}{有限型}\Term{group-extension}{扩张}。
\end{enumerate}
最值得注意的是 $K=K'$，且 $K$ 为 $k$ 的\Term{finite}{有限}不\Term{separable-extension}{可分扩张}的情形。

以 $\widehat C_{\Lambda,K}$ 表示如下\Term{category}{范畴}：\Term{object}{对象}为\Term{complete}{完备} \Term{noetherian-local-ring}{Noether 局部环} $R$，赋予 $\Lambda$-\Term{algebraic-structure}{代数结构}及 $\Lambda$-\Term{epimorphism}{满同态} $s:R\to K$。以 $C_{\Lambda,K}$ 表示其中由满足 $\operatorname{Ker}(s)$ 是\Term{finite-length}{有限长度} $R$-模的 $(R,s)$ 构成的\Term{subcategory}{子范畴}。\Term{category}{范畴} $\widehat C_{\Lambda,K}$ 可等同于 $\operatorname{Pro}(C_{\Lambda,K})$ 的一个\Term{full-subcategory}{全子范畴}。

若 $K$ 是 Artin 环，则 $C_{\Lambda,K}$ 的任意\Term{object}{对象}亦为 Artin 环。

对任意\Term{finite-dimensional}{有限维} $K'$-\Term{vector-space}{向量空间} $V$，以 $D_V(K)$ 表示由 $V$ 定义的 $K$ 上\Term{dual-number-algebra}{对偶数代数}；其元素形如 $k+v\epsilon$，其中 $k\in K$、$v\in V$、$\epsilon^2=0$，且
\[(k+v\epsilon)(k'+v'\epsilon)=kk'+(kv'+k'v)\epsilon.\]
映射 $s:D_V(K)\to K$、$s(k+v\epsilon)=k$ 使 $D_V(K)$ 成为 $C_{\Lambda,K}$ 的\Term{object}{对象}。先前 $k[\epsilon]$ 的角色在此由 $D_{K'}(K)$ 承担。

如 1.8 所述，$C_{\Lambda,K}$ 上的任意\Term{cofibered-groupoid}{余纤维群胚} $\mathcal A$ 自然延拓为 $\widehat C_{\Lambda,K}$ 上的\Term{cofibered-category}{余纤维范畴} $\widehat{\mathcal A}$。$\mathcal A$ 或 $\widehat{\mathcal A}$ 中的\Term{morphism}{态射}，若所对应的 $C_{\Lambda,K}$ 或 $\widehat C_{\Lambda,K}$ 中的\Term{algebra-homomorphism}{代数同态}\Term{surjectivity}{满射}，就称为\Term{surjectivity}{满射}。
\SourcePage{55}
假设 $\mathcal A(K)=\{e\}$。
\begin{SGAresult}{定义}{1.16}
若 $C_{\Lambda,K}$ 上的\Term{cofibered-groupoid}{余纤维群胚} $\mathcal A$ 满足下述性质，则称其为\Term{semi-homogeneous}{半齐性}的：
\begin{enumerate}
\item[(I)] $\mathcal A$ 中的每个实线图，其中 $a'\to a$ \Term{surjectivity}{满射}，都能补上虚线\Term{arrow}{箭头}成为\Term{commutative-diagram}{交换图}：
\[\begin{tikzcd}a''\arrow[r,dashed]\arrow[d,dashed]&a'\arrow[d]\\a_1\arrow[r]&a.\end{tikzcd}\]
\item[(II)] 对任意 $R$，自然映射
\[[\mathcal A(R\times_K D_{K'}(K))]\longrightarrow
[\mathcal A(R)\times\mathcal A(D_{K'}(K))]
\simeq[\mathcal A(R)]\times[\mathcal A(D_{K'}(K))]\]
为\Term{bijection}{双射}。
\end{enumerate}
若 $\mathcal A$ \Term{semi-homogeneous}{半齐性}，则 $[\mathcal A]$ 亦如此。
\end{SGAresult}
\textbf{1.17。} 先考虑 $\mathcal A$ 在 $C_{\Lambda,K}$ 中由所有 $D_V(K)$ 构成的\Term{full-subcategory}{全子范畴}上的\Term{restriction}{限制}。令 $\Omega$ 为\Term{vector-space}{向量空间} $\Omega^1_{K/\Lambda}\otimes_K K'$，$d:K\to\Omega$ 为\Term{differential}{微分}。\Term{morphism}{态射} $f:D_V(K)\to D_W(K)$ 可写成
\[f(k+v\epsilon)=k+(u(v)+D\,dk)\epsilon,\]
其中 $u\in\operatorname{Hom}_{K'}(V,W)$、$D\in\operatorname{Hom}_{K'}(\Omega,W)$。若将 $f$ 等同于 $(u,D)$，则复合法则为
\[(u,D)\circ(u',D')=(uu',D+uD').\]
条件 (II) 蕴含：

(II)$'$：从 $K'$-\Term{vector-space}{向量空间}到\Term{set}{集合}的\Term{functor}{函子}
\[V\longmapsto[\mathcal A(D_V(K))]\]
与积交换。
\begin{SGAresult}{\Term{lemma}{引理}}{1.18}
假设 $\mathcal A$ 满足条件 (II)$'$。则：
\SourcePage{56}
\begin{enumerate}
\item[(1)] $[\mathcal A(D_{K'}(K))]$ 具有唯一\Term{vector-space}{向量空间}结构，使得作为 $V$ 的\Term{functor}{函子}，
\[[\mathcal A(D_V(K))]\simeq V\otimes_{K'}[\mathcal A(D_{K'}(K))].\]
\item[(2)] 存在\Term{linear}{线性}映射
\[\gamma:\operatorname{Hom}(\Omega,K')\longrightarrow[\mathcal A(D_{K'}(K))],\]
使对任意 $D\in\operatorname{Hom}(\Omega,K')$，映射
\[[\mathcal A((k,D))]:[\mathcal A(D_{K'}(K))]\longrightarrow[\mathcal A(D_{K'}(K))]\]
为 $x\mapsto x+\gamma(D)$。
\item[(3)] 对任意 $(u,D):D_V(K)\to D_W(K)$，映射
\[(u,D):[\mathcal A(D_V(K))]\longrightarrow[\mathcal A(D_W(K))],\]
亦即
\[(u,D):V\otimes[\mathcal A(D_{K'}(K))]\longrightarrow W\otimes[\mathcal A(D_{K'}(K))],\]
为 $x\mapsto u(x)+\gamma(D)$。
\end{enumerate}
\end{SGAresult}
\begin{proof}
引入记号
\[\begin{aligned}
F(V)&=[\mathcal A(D_V(K))],\\
A&=F(K'),\\
F(u,D):F(V)\to F(W)&=\text{由 }(u,D):D_V(K)\to D_W(K)\text{ 诱导的映射},\\
F(D):F(V)\to F(V)&,\qquad F(D)=F(1,D).
\end{aligned}\]
由 (II)$'$，以及 $K'$ 是 $K'$-\Term{vector-space}{向量空间}\Term{category}{范畴}中的 $K'$-\Term{vector-space}{向量空间}\Term{object}{对象}这一事实，$F(K')$ 是 $(\mathrm{Ens})$ 中的 $K'$-\Term{vector-space}{向量空间}，因而是 $K'$-\Term{vector-space}{向量空间}；(1) 即显然。现等同 $F(V)$ 与 $V\otimes A$。

由恒等式 $(u,D)=(1,D)\circ(u,0)$ 得
\[F(u,D)=F(D)\circ(u\otimes1_A).\]
$F$ 的\Term{functorial}{函子性}即为
\[\begin{aligned}
F(0)&=\operatorname{Id},\\
F(D)\circ F(D')&=F(D+D'),\\
(u\otimes1)\circ F(D)&=F(uD)\circ(u\otimes1_A).
\end{aligned}\]
对任意 $D:\Omega\to V$ 及 $n\in\mathbf N$，若 $i,p$ 为映射
\[i:V\hookrightarrow V\oplus K'^n,\qquad p:V\oplus K'^n\longrightarrow K'^n,\]
\SourcePage{57}
则 $piD=0$，且下图交换：
\[\begin{tikzcd}
V\otimes A\arrow[r]\arrow[d,"F(D)"']&(V\oplus K'^n)\otimes A\arrow[r]\arrow[d,"{F(i,D)}"]&K'^n\otimes A\arrow[d,equal]\\
V\otimes A\arrow[r]&(V\oplus K'^n)\otimes A\arrow[r]&K'^n\otimes A.
\end{tikzcd}\]
设 $x:K^n\to V$ 为映射，$x_i=x(e_i)$，$a_i$ 为 $A$ 中的一族元素。有 $(1_V+x)\circ i\circ D=D$。因此，若
\[F(iD)\Bigl(\sum e_i a_i\Bigr)=\sum e_i a_i+a\qquad(a\in V\otimes H),\]
则
\[F(D)\Bigl(\sum x_i a_i\Bigr)=\sum x_i a_i+a,\]
且 $a$ 不依赖于各 $x_i$。所以
\[F(D)(y)=y+G(D),\qquad\text{其中}\qquad
\begin{cases}G(D+D')=G(D)+G(D'),\\uG(D)=F(uD)\quad(\text{对任意 }u:V\to W).\end{cases}\]
容易验证，每个这样的 $G$ 都由某个 $\gamma\in\Omega\otimes A$ 经公式
\[G(D)=D\otimes1_A(\gamma)\]
定义，这正是 (2)、(3) 所断言的。
\end{proof}
\begin{SGAresult}{定义}{1.19}
(1) 若 $\widehat{\mathcal A}$ 中的\Term{object}{对象} $\xi$ 满足：$\mathcal A$ 中每个\Term{surjectivity}{满射} $\alpha:a'\to a$ 均诱导\Term{surjectivity}{满射}
\[\operatorname{Hom}_{\widehat{\mathcal A}}(\xi,a')\longrightarrow
\operatorname{Hom}_{\widehat{\mathcal A}}(\xi,a),\]
则称 $\xi$ 为\Term{versal-formal-object}{通用形式对象}。

(2) 设 $\xi\in\widehat{\mathcal A}$ 为\Term{versal-formal-object}{通用形式对象}，其在 $\widehat C_{\Lambda,K}$ 中的像为 $(R,s)$。若对任意两个映射 $\varphi_1,\varphi_2:R\rightrightarrows D_{K'}(K)$，由 $\varphi_1(\xi)\simeq\varphi_2(\xi)$ 可\Term{pushout}{推出}：存在 $K$ 到 $K'$ 的\Term{derivation}{导子} $D$，使
\[(1,D)(\varphi_1)=\varphi_2,\qquad\text{即 }\varphi_2-\varphi_2=Ds,\]
则称 $\xi$ 为极小的。
\end{SGAresult}
\begin{SGAresult}{\Term{theorem}{定理}}{1.20}
若 $\mathcal A$ \Term{semi-homogeneous}{半齐性}，且
\[\dim_{K'}[\mathcal A(D_{K'}(K))]<\infty,\]
则 $\mathcal A$ 存在\Term{minimally-versal-formal-object}{极小通用形式对象}，且任意两个这样的\Term{object}{对象}\Term{isomorphism}{同构}。
\SourcePage{58}
\end{SGAresult}
\begin{proof}
沿用 1.18 的记号，设 $H^*$ 是 $[\mathcal A(D_{K'}(K))]$ 中与 $\gamma$ 的像互补的\Term{vector-subspace}{子空间}，$H$ 为 $H^*$ 的\Term{duality}{对偶}。定义 $R_0=D_H(K)$。取 $\mathcal A(R_0)$ 中的\Term{object}{对象} $\xi_0$，使其在
\[[\mathcal A(R_0)]\simeq H\otimes[\mathcal A(D_{K'}(K))]
\simeq\operatorname{Hom}(H^*,[\mathcal A(D_{K'}(K))])\]
中的像为 $H^*$ 到 $[\mathcal A(D_{K'}(K))]$ 的\Term{inclusion}{包含映射}。于是：
\begin{enumerate}
\item[(1)] 映射 $\operatorname{Hom}_{C_{\Lambda,K}}(D_H(K),D_V(K))\to[\mathcal A(D_V(K))]$ \Term{surjectivity}{满射}；
\item[(2)] 极小性条件 1.19(2) 对 $(R_0,\xi_0)$ 成立。
\end{enumerate}
现选取 $S\in\widehat C_{\Lambda,K}$，使其在 $\Lambda$ 上\Term{formally-smooth}{形式光滑}，并取\Term{epimorphism}{满同态} $t:S\to R_0$，使
\[\Omega_{R_0/\Lambda}\otimes K'\xrightarrow\sim\Omega_{S/\Lambda}\otimes K'.\]
若 $M$ 为 $S$ 的\Term{maximal-ideal}{极大理想}，则对 $n\geq1$，由下述性质归纳定义 $R_n=S/I_n$ 及 $\xi_n\in\operatorname{ob}(\mathcal A(R_n))$：
\begin{enumerate}
\item[(a)] $I_{n+1}$ 是所有满足 $I_nM\subset I\subset I_n$、且 $\xi_n$ 能延拓为 $\mathcal A(S/I)$ 中\Term{object}{对象} $\xi$ 的\Term{ideal}{理想} $I$ 的交；
\item[(b)] $\xi_{n+1}$ 是 $\xi_n$ 的一个延拓。
\end{enumerate}
若 $R=\varprojlim R_n$、$\xi=\varprojlim\xi_n$，则可验证 $(R,\xi)$ 为\Term{minimally-versal-formal-object}{极小通用形式对象}，并且在\Term{isomorphism}{同构}意义下唯一。
\end{proof}
